% ================================================================
\section{Complementos}
% ================================================================

\subsection{Tabelas hash (pesquisa)}\label{sec:hash}
\begin{complemento}[ -- pesquisa proposta em 02/09, entrega 19/09 (Cormen, cap.~11; Szwarcfiter, cap.~10)]
\textbf{Motivação:} implementar um \emph{dicionário} (inserir, buscar, remover por chave) com acesso ``direto''. Endereçamento direto (vetor indexado pela chave) é $O(1)$ mas exige memória do tamanho do universo de chaves (infinito ou enorme). Uma \textbf{função hash} $h:U\to\{0,\dots,m-1\}$ transforma a chave num endereço de uma tabela de $m$ posições.

\textbf{Funções clássicas:} \emph{divisão} $h(k)=k\bmod m$ ($m$ primo longe de potências de 2); \emph{multiplicação} $h(k)=\floor{m\,(kA\bmod1)}$ ($A\approx(\sqrt5-1)/2$); \emph{dobra}/\emph{meio do quadrado}; para strings, polinomial $h(s)=\sum s_i\,b^{i}\bmod m$; \emph{hashing universal} (escolher $h$ aleatoriamente de uma família, p.ex.\ $((ak+b)\bmod p)\bmod m$), que garante bom caso esperado para qualquer entrada.

\textbf{Colisões} ($h(k_1)=h(k_2)$, $k_1\ne k_2$) são inevitáveis (princípio da casa dos pombos: $|U|>m$). Tratamento:
\begin{itemize}
  \item \textbf{Encadeamento exterior:} cada posição aponta para uma lista das chaves com aquele endereço.
  \item \textbf{Encadeamento interior:} as listas usam posições da própria tabela (uma zona de colisões, ou ponteiros entre posições livres).
  \item \textbf{Endereçamento aberto:} tudo fica na tabela; em colisão, sonda-se uma sequência de posições: \emph{linear} $h(k)+i$ (agrupamento primário), \emph{quadrática} $h(k)+c_1i+c_2i^2$, \emph{duplo hashing} $h_1(k)+i\,h_2(k)$. Remoção exige marca ``removido''.
\end{itemize}
\textbf{Complexidade:} com fator de carga $\alpha=n/m$ e hash uniforme, busca em $\Theta(1+\alpha)$ esperado com encadeamento, e $\le\frac1{1-\alpha}$ sondagens esperadas (busca sem sucesso) no endereçamento aberto; mantendo $\alpha=O(1)$ (redimensionando a tabela, custo amortizado $O(1)$), as operações são $O(1)$ esperado. \textbf{Pior caso: $\Theta(n)$} (todas as chaves no mesmo endereço). Hash não mantém ordem (sem mínimo, sucessor, intervalos --- aí usa-se AVL/rubro-negra).
\end{complemento}

\subsection{Programação dinâmica}\label{sec:pd}
\begin{complemento}[ -- slides de PD (subsequência comum máxima); não previsto no formato da AV1]
Na decomposição recursiva, \textbf{subproblemas idênticos} podem aparecer muitas vezes, gerando tempo exponencial mesmo com poucos subproblemas distintos (ex.: Fibonacci recursivo). \textbf{PD}: guardar a solução de cada subproblema numa tabela e resolvê-los de baixo para cima (\emph{bottom-up}) ou com memorização. D\&C é bom quando os subproblemas são disjuntos; PD quando eles se sobrepõem.

\textbf{Subsequência comum máxima (LCS).} $X=x_1..x_n$, $Y=y_1..y_m$; $C(i,j)=$ tamanho da LCS de $X_i$ e $Y_j$:
\[
C(i,j)=\begin{cases}0 & i=0\text{ ou }j=0\\ 1+C(i-1,j-1) & x_i=y_j\\ \max\{C(i-1,j),C(i,j-1)\} & x_i\ne y_j\end{cases}
\]
Preenche-se a tabela linha a linha: $O(nm)$ tempo e espaço; resposta $C(n,m)$. Ex.: $X=5,2,1,3,9,10,14$, $Y=2,3,1,10$: LCS de tamanho 3 ($2,3,10$ ou $2,1,10$). Para \emph{recuperar} a subsequência, volte de $(n,m)$: se $x_i=y_j$, $x_i$ faz parte e vá a $(i-1,j-1)$; senão vá para a vizinha de maior valor.
\end{complemento}
